\documentclass[12pt,a4paper]{article}
\usepackage[utf8]{inputenc}
\usepackage[T1]{fontenc}
\usepackage{geometry}
\geometry{margin=2.5cm}
\usepackage{hyperref}
\usepackage{booktabs}
\usepackage{enumitem}
\usepackage{listings}
\usepackage{xcolor}

\lstset{
  basicstyle=\ttfamily\small,
  breaklines=true,
  frame=single,
  backgroundcolor=\color{gray!10}
}

\title{ShinySwarm --- Repository Context Document}
\author{Iva Gunjaca \\ University of Amsterdam \\ MSc Software Engineering (18 ECTS)}
\date{Last updated: June 2026}

\begin{document}
\maketitle

\tableofcontents
\newpage

% ==================================================================
\section{Project Overview}
% ==================================================================

ShinySwarm is a collaborative R~Shiny platform built as a Master's thesis
project at the University of Amsterdam (UvA), Software Engineering track.
It enables real-time collaboration on R~Shiny applications --- multiple users
can simultaneously interact with the same Shiny app, see each other's changes,
save/restore checkpoints, and share analytical results. The platform is designed
for integration with NaaVRE (Notebook-as-a-Virtual-Research-Environment) at UvA.

% ==================================================================
\section{Repository Structure}
% ==================================================================

The repository contains three architectures for comparison:

\begin{description}[style=nextline]
  \item[\texttt{Thesis-Project-Final-RESTAPI/}]
    REST architecture. Spring Boot backend (Java~21, Spring Security, JPA/Hibernate,
    PostgreSQL). Angular~19 frontend. Plumber R backends for each Shiny app.
    State synchronisation via Redis polling
    (POST state $\rightarrow$ Redis $\rightarrow$ GET poll).

  \item[\texttt{Thesis-Project-Final-Kafka/}]
    Kafka architecture. Same Spring Boot + Angular stack, but state synchronisation
    via Apache Kafka topics (produce to \texttt{input} topic $\rightarrow$ Kafka
    consumer $\rightarrow$ \texttt{output} topic $\rightarrow$ Plumber polls).
    Confluent Kafka + Zookeeper.

  \item[\texttt{whole\_apps/}]
    Monolithic baseline. 8 standalone Shiny apps with no collaboration layer.
    Docker Compose + Nginx reverse proxy. Used as the computation floor for
    benchmarking.
\end{description}

% ==================================================================
\section{The 8 Shiny Applications}
% ==================================================================

Each architecture runs the same 8 applications:

\begin{enumerate}[nosep]
  \item \textbf{Calculator} --- simple arithmetic (instant computation)
  \item \textbf{Visual Analytics} --- mtcars dataset filtering with checkboxes (cyl values)
  \item \textbf{Advanced Visual Analytics} --- airquality dataset with month/variable filters
  \item \textbf{Data Exchange} --- CSV upload, processing (uppercase transformation), download
  \item \textbf{Climate Anomaly Detector} --- sensor CSV upload, anomaly detection, date parsing
  \item \textbf{Geospatial Editor} --- Leaflet map with marker placement (delta sync)
  \item \textbf{Monte Carlo Simulator} --- population viability analysis (15--30s computation)
  \item \textbf{ML Trainer (Habitat Suitability AI)} --- Random Forest training with ranger (30--60s computation)
\end{enumerate}

% ==================================================================
\section{Technology Stack}
% ==================================================================

\begin{table}[h]
\centering
\begin{tabular}{@{}ll@{}}
\toprule
\textbf{Layer} & \textbf{Technology} \\
\midrule
Backend        & Spring Boot 3.5, Java 21, Spring Security (JWT), Spring Data JPA, PostgreSQL 15 \\
State sync     & Redis (REST architecture) / Kafka 3.9 + Zookeeper (Kafka architecture) \\
Frontend       & Angular 19, TypeScript, Tailwind CSS, STOMP-over-WebSocket (presence) \\
R backends     & Plumber API servers (one per Shiny app), state serialisation as JSON \\
Shiny frontends & R Shiny apps embedded in iframes, HTTP polling to Plumber backends \\
Infrastructure & Kubernetes (k3s), Hetzner CX32 (8 vCPU, 16\,GB RAM), Docker, Rancher \\
Testing        & Playwright (end-to-end, 12 test files A--L), k6 (load testing, 10 benchmarks) \\
CI/CD          & GitHub Actions (build Docker images on push to main, push to GHCR) \\
\bottomrule
\end{tabular}
\end{table}

% ==================================================================
\section{Key Architectural Patterns}
% ==================================================================

\begin{itemize}[nosep]
  \item \textbf{Full-state serialisation:} Both inputs AND outputs are captured
        as JSON (not just inputs like \texttt{shinyURL}/bookmarking). This enables
        exact reproducibility for non-deterministic analyses.

  \item \textbf{Collaboration modes:} Solo, real-time sync, role-based access
        (OWNER/EDITOR/VIEWER), checkpoint save/restore, cross-user checkpoint
        sharing (RQ5).

  \item \textbf{Cross-user checkpoint sharing:} The \texttt{SavedState} entity
        has a \texttt{sessionId} column. Any session participant can list and
        restore checkpoints saved by other participants. Authorisation: owner
        OR session participant.

  \item \textbf{Solo-mode guard:} Session IDs starting with \texttt{"solo-"}
        are filtered out to prevent accidental session association.
\end{itemize}

% ==================================================================
\section{Research Questions}
% ==================================================================

\begin{quote}
\textbf{Central RQ:} \textit{How can collaborative, containerised R~Shiny
applications be architected to support reproducible, real-time multi-user
scientific analyses?}
\end{quote}

\begin{description}[style=nextline, nosep]
  \item[RQ1] What collaboration modes and permission models are needed for
             interactive Shiny applications?
  \item[RQ2] How can application state (inputs + outputs) be shared between
             users without recomputation?
  \item[RQ3] What are the trade-offs between REST (synchronous) and Kafka
             (asynchronous) state synchronisation?
  \item[RQ4] How do the architectures compare in latency, throughput, data loss,
             and cross-contamination?
  \item[RQ5] Does containerised deployment + checkpoint-based state persistence
             improve the reproducibility of collaborative R~Shiny analyses?
\end{description}

% ==================================================================
\section{Benchmarking (k6)}
% ==================================================================

\subsection{Test Suite}

10 benchmarks, each targeting a specific performance dimension:

\begin{table}[h]
\centering
\small
\begin{tabular}{@{}cll@{}}
\toprule
\textbf{\#} & \textbf{Test} & \textbf{Measures} \\
\midrule
00 & Baseline           & Per-app response time (1 VU, 5 reps) \\
01 & State relay        & POST $\rightarrow$ GET round-trip latency (20 VUs) \\
02 & Collaboration      & State propagation between users (30 VUs) \\
03 & Save/restore       & Checkpoint persistence latency (15 VUs) \\
04 & Throughput          & Requests/sec under load (100 VUs) \\
05 & Data loss          & Percentage of states lost in transit \\
06 & Cross-contamination & State leaking between sessions \\
07 & Session lifecycle   & Create/join/leave overhead (50 VUs) \\
08 & Multi-user collab   & 3 VUs collaborating on each app \\
09 & WebSocket presence  & STOMP connect + roundtrip (50 VUs) \\
\bottomrule
\end{tabular}
\end{table}

\subsection{Threshold Calibration}

Per-app baseline-calibrated thresholds in \texttt{k6/config.js}. Each app has
its own threshold based on baseline median $\times$ multiplier:

\begin{itemize}[nosep]
  \item $10\times$ for relay, propagation, lifecycle (single-threaded queuing)
  \item $50\times$ for throughput (100 VUs)
  \item $10\times$ for save/restore (PostgreSQL)
\end{itemize}

\subsection{Execution Plan}

\begin{itemize}[nosep]
  \item 5 runs per architecture $\times$ 2 locations (co-located on server +
        remote from laptop) for statistical validity
  \item Report mean $\pm$ 95\% CI ($\text{mean} \pm 1.96 \times \sigma / \sqrt{n}$, $n=5$)
  \item Automated via \texttt{k6/run-suite.sh} (cooldown between tests, sleep
        prevention, numbered output folders)
\end{itemize}

% ==================================================================
\section{Playwright Tests}
% ==================================================================

12 test files ordered A--L, running sequentially (\texttt{fullyParallel: false}):

\begin{table}[h]
\centering
\small
\begin{tabular}{@{}cll@{}}
\toprule
\textbf{File} & \textbf{Name} & \textbf{Scope} \\
\midrule
A & Login basics            & Login page loads, auth redirect \\
B & Handshake               & WebSocket connection \\
C & Collab sync             & Alice \& Bob see each other \\
D & Cross-user restore      & RQ5 standalone (Calculator) + isolation test \\
E & Calculator              & 5 tests: solo, collab, security, restore, RQ5 \\
F & Visual Analytics        & 5 tests: solo, collab, security, restore, RQ5 \\
G & Advanced Analytics      & 5 tests: solo, collab, security, restore, RQ5 \\
H & Data Exchange           & 5 tests: solo, collab, security, restore, RQ5 \\
I & Climate Anomaly         & 5 tests: solo, collab, security, restore, RQ5 \\
J & Geospatial Delta        & 6 tests: solo, collab, security, restore, RQ5, delta sync \\
K & Monte Carlo             & 5 tests: solo, collab, security, restore, RQ5 \\
L & ML Trainer              & 5 tests: solo, collab, security, restore, RQ5 \\
\bottomrule
\end{tabular}
\end{table}

\subsection{Per-App Test Pattern (Tests 1--5)}

\begin{enumerate}[nosep]
  \item \textbf{Solo mode:} Compute \& save state
  \item \textbf{Collab mode:} Real-time synchronisation (Alice sends, Bob receives)
  \item \textbf{Security:} Role-based permission (demote to VIEWER, verify inputs disabled)
  \item \textbf{Time machine:} Restore checkpoint in solo mode
  \item \textbf{RQ5 cross-user restore:} Alice saves, changes state to something
        distinguishable, leaves session. Bob joins, restores Alice's checkpoint,
        verifies he gets the saved state (not Alice's later changes)
\end{enumerate}

% ==================================================================
\section{Key Files}
% ==================================================================

\begin{table}[h]
\centering
\small
\begin{tabular}{@{}ll@{}}
\toprule
\textbf{File} & \textbf{Purpose} \\
\midrule
\texttt{k6/config.js}                    & App configs, thresholds, BASE\_URL \\
\texttt{k6/run-suite.sh}                 & Automated benchmark runner (N runs) \\
\texttt{k6/00-baseline.js} -- \texttt{09-*.js} & Individual k6 benchmark scripts \\
\texttt{docs/thesis/introduction\_and\_research\_questions.tex} & RQs, methodology \\
\texttt{docs/thesis/literature\_review\_summary.tex}            & 30+ sources \\
\texttt{docs/thesis/architecture\_summary.tex}                  & Full architecture docs \\
\texttt{docs/thesis/next-steps-benchmarking.md}                 & Remaining work roadmap \\
\texttt{docs/thesis/repository-context.tex}                     & This document \\
\texttt{k6/results/kafka/thesis-notes.tex}                      & Working thesis notes \\
\texttt{k8s/deploy-rest.sh}              & Deploy REST to Kubernetes \\
\texttt{k8s/deploy-kafka.sh}             & Deploy Kafka to Kubernetes \\
\texttt{k8s/deploy-monolithic.sh}        & Deploy monolithic to Kubernetes \\
\texttt{k8s/teardown.sh}                 & Tear down current deployment \\
\texttt{.github/workflows/build-images.yml} & CI: build \& push Docker images \\
\bottomrule
\end{tabular}
\end{table}

% ==================================================================
\section{Infrastructure and Deployment}
% ==================================================================

\begin{itemize}[nosep]
  \item \textbf{Cluster:} Hetzner CX32, single node, 8 vCPU, 16\,GB RAM
  \item \textbf{Server IP:} 188.245.60.172
  \item \textbf{REST NodePort:} 30001
  \item \textbf{Kafka NodePort:} 30002
  \item \textbf{Test users:} alice, bob, charlie (password: \texttt{password})
  \item \textbf{Docker images:} \texttt{ghcr.io/ifapifa/collaborative-rshiny/*}
  \item \textbf{Kubernetes:} k3s, managed via Rancher
\end{itemize}

% ==================================================================
\section{Conventions and Preferences}
% ==================================================================

\begin{itemize}[nosep]
  \item Never push or commit without explicit permission from the student.
  \item Explain complex concepts simply when asked.
  \item The thesis is due in approximately 4 months (October 2026).
  \item All three architectures (REST, Kafka, monolithic) must be benchmarked
        with the same test suite for fair comparison.
  \item Playwright tests are identical across REST and Kafka architectures
        (same test files copied to both).
  \item The \texttt{DataInitializer} seeds test users on Spring Boot startup
        if they do not already exist in PostgreSQL.
\end{itemize}

% ==================================================================
\section{Current Status (June 2026)}
% ==================================================================

\begin{itemize}[nosep]
  \item[\checkmark] REST architecture deployed and benchmarked (1 run, all 10 tests)
  \item[\checkmark] Per-app baseline-calibrated thresholds implemented
  \item[\checkmark] Cross-user checkpoint sharing implemented (both architectures)
  \item[\checkmark] RQ5 Playwright tests added to all 8 apps (both architectures)
  \item[\checkmark] Test files renamed A--L for ordered execution
  \item[\checkmark] \texttt{run-suite.sh} script for automated 5$\times$ runs
  \item[$\circ$] Kafka benchmarks (in progress --- baseline running)
  \item[$\circ$] REST benchmarks re-run 5$\times$ for statistical validity
  \item[$\circ$] Monolithic benchmarks (deploy, instrument with \texttt{system.time()}, run)
  \item[$\circ$] Data analysis: parse results, compute mean/CI, generate plots
  \item[$\circ$] Thesis writing: Results, Discussion, Threats to Validity, Conclusion
\end{itemize}

\end{document}
